\documentclass[11pt]{article}
\usepackage[margin=1in]{geometry}
\usepackage{amsmath,amssymb,amsthm}
\usepackage{graphicx}
\usepackage{float}
\usepackage[T1]{fontenc}
\usepackage{lmodern}
\usepackage{microtype}
\usepackage[hidelinks]{hyperref}

\newtheorem{theorem}{Theorem}
\newtheorem{lemma}[theorem]{Lemma}
\newtheorem{proposition}[theorem]{Proposition}
\newcommand{\norm}[1]{\lVert#1\rVert}
\newcommand{\abs}[1]{\lvert#1\rvert}
\newcommand{\Bcal}{\mathcal B}
\newcommand{\Pinf}{\mathcal P_\infty}
\newcommand{\WOH}{\widehat{\mathrm{WOH}}}

\title{Finite Scale Meshes Give the Exact Borel Complexity\\
of Weak Octahedrality}
\author{Full solution packet for arXiv:2404.19457}
\date{22 July 2026}

\begin{document}
\maketitle

\begin{center}
\fbox{\parbox{0.91\textwidth}{\textbf{Status: full solution likely valid.}
In each of the Polish coding spaces $\Pinf$ and $\Bcal$, the isometry class
$\WOH$ of weakly octahedral separable Banach spaces is
$G_\delta$-complete.  This improves the source paper's
$F_{\sigma\delta}$ upper bound and answers its exact-complexity question in
both coding spaces.}}
\end{center}

\section{Source question}

L\'opez-P\'erez, Mart\'inez Va\~n\'o, and Rueda Zoca prove that $\WOH$ is
$F_{\sigma\delta}$ in $\Pinf$ and $\Bcal$ \cite[Theorem 6.6]{source}.
On source PDF pages 24--25 they ask for the exact complexity of weakly
octahedral spaces.  They explain that a positive answer to the related
single-scale question would give a $G_\delta$ description in $\Bcal$.

More precisely, the single-scale question asks whether weak octahedrality is
equivalent to the following condition: for every $\varepsilon>0$, every
$x_1,\ldots,x_n\in S_X$, and every $f\in B_{X^*}$, there is $y\in S_X$ such
that
\begin{equation}\label{eq:single}
 \norm{x_i+y}\geq (1-\varepsilon)(\abs{f(x_i)}+1)
 \qquad (i=1,\ldots,n).
\end{equation}
That characterization remains outside the scope of this packet.  We instead
settle the exact-complexity question directly.

\begin{figure}[H]
 \centering
 \includegraphics[width=0.98\textwidth]{figures/open_problem_crop.png}
 \caption{The open questions in arXiv:2404.19457, source PDF pages 24--25.}
\end{figure}

\section{Definitions and notation}

All spaces are real.  Recall the finite-set characterization
\cite[Proposition 2.2]{HLP}: $X$ is \emph{weakly octahedral} if and only if,
for every $\eta>0$, $x_1,\ldots,x_n\in S_X$, and $f\in B_{X^*}$, there is
$y\in S_X$ such that
\begin{equation}\label{eq:woh}
 \norm{x_i+t y}\geq (1-\eta)(\abs{f(x_i)}+t)
 \quad (i=1,\ldots,n,\ t>0).
\end{equation}

We use the source paper's coding.  Let $V$ be the countable rational vector
space of finitely supported sequences and let $\mu$ be a seminorm code.
The quotient completion is denoted by $X_\mu$.  The spaces $\Pinf$ and
$\Bcal$ consist, respectively, of the codes for infinite-dimensional spaces
and the codes for which $\mu$ is a norm on $V$; see \cite{CDDK,source}.
For $v\in V$ with $\mu(v)>0$, put
\[
 x_v^\mu=\frac{[v]}{\mu(v)}\in S_{X_\mu}.
\]
The set of all such normalized rational directions is dense in $S_{X_\mu}$.

\section{Proof intuition}

There are two compressions.  First, at any requested accuracy, the continuum
of scales $t>0$ in \eqref{eq:woh} is unnecessary.  The reverse triangle
inequality handles very small and very large $t$.  On the remaining compact
interval, $t\mapsto\norm{x+t y}$ is $1$-Lipschitz, so a fixed finite rational
mesh controls every scale.

Second, for finitely many test directions we do not need a global code for
the entire dual ball.  We retain only the finite evaluation vector
$(f(x_1),\ldots,f(x_n))\in[-1,1]^n$.  The admissible evaluation vectors form
a compact set with a closed graph: membership is exactly the countable family
of Hahn--Banach inequalities for rational linear combinations.  This remains
true in $\Pinf$, where the source paper's global dual-ball graph is not
closed.  The finite-mesh witness relation is open; compactness of
$[-1,1]^n$ then makes its universal co-projection open.  A countable
intersection gives the desired $G_\delta$ definition.

\section{The finite scale lemma}

Fix an integer $m\geq2$ and write $\varepsilon=1/m$.  Set
\[
 a_\varepsilon=\frac{\varepsilon}{4},\qquad
 b_\varepsilon=\frac{2}{\varepsilon},\qquad
 h_\varepsilon=\frac{\varepsilon^2}{64}.
\]
Choose once and for all a finite rational $h_\varepsilon$-net
$T_\varepsilon$ in $[a_\varepsilon,b_\varepsilon]$.

\begin{lemma}[finite scale mesh]\label{lem:mesh}
Let $x,y\in S_X$ and $c\in[0,1]$.  If
\begin{equation}\label{eq:gridhalf}
 \norm{x+s y}\geq(1-\varepsilon/2)(c+s)
 \qquad(s\in T_\varepsilon),
\end{equation}
then
\begin{equation}\label{eq:allscale}
 \norm{x+t y}\geq(1-\varepsilon)(c+t)
 \qquad(t>0).
\end{equation}
\end{lemma}

\begin{proof}
If $0<t\leq a_\varepsilon$, then
\[
 \norm{x+t y}\geq1-t
 \geq(1-\varepsilon)(1+t)
 \geq(1-\varepsilon)(c+t).
\]
The middle inequality follows by checking the right endpoint
$t=\varepsilon/4$.  If $t\geq b_\varepsilon$, then
\[
 \norm{x+t y}\geq t-1
 \geq(1-\varepsilon)(1+t)
 \geq(1-\varepsilon)(c+t),
\]
because $\varepsilon(1+t)\geq2$.

It remains to take $t\in[a_\varepsilon,b_\varepsilon]$.  Choose
$s\in T_\varepsilon$ with $\abs{t-s}\leq h_\varepsilon$.  By
\eqref{eq:gridhalf} and the triangle inequality,
\begin{align*}
 \norm{x+t y}
 &\geq (1-\varepsilon/2)(c+s)-h_\varepsilon\\
 &\geq (1-\varepsilon/2)(c+t)-2h_\varepsilon.
\end{align*}
Since $c+t\geq a_\varepsilon=\varepsilon/4$ and
$2h_\varepsilon=\varepsilon^2/32\leq
(\varepsilon/2)(c+t)$, this is at least
$(1-\varepsilon)(c+t)$.
\end{proof}

\begin{lemma}[joint continuity of finite-dimensional evaluation]
\label{lem:continuity}
Fix a finite-dimensional coordinate subspace $E\subset c_{00}$.  If
$\mu_k\to\mu$ pointwise on $V$ and $z_k\to z$ in $E$, then
\[
 \bar\mu_k(z_k)\longrightarrow\bar\mu(z),
\]
where $\bar\mu$ denotes the real extension of the rational seminorm.
Consequently, all normalized norm expressions used below are continuous on
the open set where their denominators are positive.
\end{lemma}

\begin{proof}
Choose coordinate vectors $e_1,\ldots,e_d$ spanning $E$.  Pointwise
convergence makes $\sup_k\mu_k(e_j)<\infty$ for every $j$.  Hence the
$\bar\mu_k$ are equi-Lipschitz on $E$ for the coordinate $\ell_1$ norm.
Approximate $z$ by a rational vector $r\in E\cap V$, use pointwise
convergence at $r$, and then use the uniform Lipschitz bound to compare
$z_k,r$, and $z$.  Letting the rational approximation error tend to zero
proves the claim.
\end{proof}

\section{A robust countable criterion}

\begin{proposition}\label{prop:criterion}
A Banach space $X$ is weakly octahedral if and only if the following holds.
For every $m\geq2$, every finite family $z_1,\ldots,z_n$ from a fixed dense
subset of $S_X$, and every $f\in B_{X^*}$, there is a $y$ in that dense
subset such that
\begin{equation}\label{eq:criterion}
 \norm{z_i+s y}>
 (1-\varepsilon/4)(\abs{f(z_i)}+s)
 \quad(i=1,\ldots,n,\ s\in T_\varepsilon),
 \qquad \varepsilon=1/m.
\end{equation}
\end{proposition}

\begin{proof}
Suppose first that $X$ is weakly octahedral.  Apply \eqref{eq:woh} to the
$z_i$ with error $\varepsilon/8$, obtaining $y_0\in S_X$.  Choose a dense
point $y$ with $\norm{y-y_0}<\varepsilon/16$.  For $s>0$,
\begin{align*}
 \norm{z_i+s y}
 &\geq (1-\varepsilon/8)(\abs{f(z_i)}+s)
       -s\varepsilon/16\\
 &>(1-\varepsilon/4)(\abs{f(z_i)}+s),
\end{align*}
so \eqref{eq:criterion} holds.

Conversely, let $x_1,\ldots,x_n\in S_X$, $f\in B_{X^*}$, and a target
error $\eta>0$ be given.  Choose $m$ so large that
$\varepsilon=1/m<\eta$.  Pick dense points $z_i$ satisfying
$\norm{x_i-z_i}<\varepsilon^2/64$.  Let $y$ be supplied by
\eqref{eq:criterion}.  For $s\in T_\varepsilon$, using
$\abs{f(z_i)-f(x_i)}\leq\norm{z_i-x_i}$, we get
\begin{align*}
 \norm{x_i+s y}
 &\geq \norm{z_i+s y}-\norm{x_i-z_i}\\
 &>(1-\varepsilon/4)(\abs{f(x_i)}+s)
      -2\norm{x_i-z_i}\\
 &\geq(1-\varepsilon/2)(\abs{f(x_i)}+s).
\end{align*}
The last step uses $s\geq a_\varepsilon=\varepsilon/4$.  Lemma
\ref{lem:mesh}, with $c=\abs{f(x_i)}$, now gives
\[
 \norm{x_i+t y}\geq(1-\varepsilon)(\abs{f(x_i)}+t)
 >(1-\eta)(\abs{f(x_i)}+t)
\]
for every $i$ and $t>0$.  This is \eqref{eq:woh}.
\end{proof}

\section{Closed finite traces of the dual ball}

Fix nonzero $v_1,\ldots,v_n\in V$ and let
\[
 D_{\mathbf v}=\{\mu:\mu(v_i)>0\text{ for }i=1,\ldots,n\}.
\]
This is open in both $\Pinf$ and $\Bcal$ (and equals all of $\Bcal$ for a
nonzero tuple).  For $\mu\in D_{\mathbf v}$ define
\[
 C_{\mu,\mathbf v}
 =\{(f(x_{v_1}^\mu),\ldots,f(x_{v_n}^\mu)):
       f\in B_{X_\mu^*}\}\subset[-1,1]^n.
\]

\begin{lemma}\label{lem:trace}
The set
\[
 \mathfrak C_{\mathbf v}
 =\{(\mu,\alpha)\in D_{\mathbf v}\times[-1,1]^n:
                 \alpha\in C_{\mu,\mathbf v}\}
\]
is relatively closed.  Equivalently, $\alpha\in C_{\mu,\mathbf v}$ if
and only if, for every $q=(q_1,\ldots,q_n)\in\mathbb Q^n$,
\begin{equation}\label{eq:HB}
 \abs{\sum_{i=1}^n q_i\alpha_i}
 \leq
 \bar\mu\!\left(\sum_{i=1}^n
             \frac{q_i v_i}{\mu(v_i)}\right).
\end{equation}
\end{lemma}

\begin{proof}
Every evaluation vector of a norm-one functional satisfies
\eqref{eq:HB}.  Conversely, the inequalities show that the assignment
$x_{v_i}^\mu\mapsto\alpha_i$ induces a well-defined linear functional of
norm at most one on their span: norm-zero rational relations are respected,
and real relations follow by density.  Hahn--Banach extends it to
$B_{X_\mu^*}$.  Finally, Lemma \ref{lem:continuity} makes the right side of
\eqref{eq:HB} continuous on $D_{\mathbf v}$.  Thus the countable
intersection of the closed inequalities \eqref{eq:HB} is relatively closed.
\end{proof}

\section{\texorpdfstring{The $G_\delta$ description}{The G-delta description}}

Fix $m,n$ and $\mathbf v=(v_1,\ldots,v_n)$.  On
$D_{\mathbf v}\times[-1,1]^n$, let $R_{m,n,\mathbf v}$ consist of the pairs
$(\mu,\alpha)$ for which either $\alpha\notin C_{\mu,\mathbf v}$, or there
is a $w\in V$ with $\mu(w)>0$ such that, for every $i$ and every
$s\in T_\varepsilon$,
\begin{equation}\label{eq:openwitness}
 \bar\mu\!\left(
       \frac{v_i}{\mu(v_i)}+s\frac{w}{\mu(w)}\right)
 >(1-\varepsilon/4)(\abs{\alpha_i}+s).
\end{equation}
By Lemmas \ref{lem:continuity} and \ref{lem:trace}, this relation is open:
the inadmissibility alternative is open, and the witness alternative is a
countable union over $w$ of finite intersections of strict continuous
inequalities.

Put
\[
 O_{m,n,\mathbf v}
 =\{\mu\in D_{\mathbf v}:
       \{\mu\}\times[-1,1]^n\subset R_{m,n,\mathbf v}\}.
\]
Since the cube is compact and $R_{m,n,\mathbf v}$ is open,
$O_{m,n,\mathbf v}$ is open relative to $D_{\mathbf v}$ (the tube lemma).
Therefore
\[
 S_{m,n,\mathbf v}=D_{\mathbf v}^{,c}\cup O_{m,n,\mathbf v}
\]
is $G_\delta$: in a metric space the union of a closed set and an open set
is $G_\delta$.

\begin{theorem}\label{thm:main}
In both $\Pinf$ and $\Bcal$,
\begin{equation}\label{eq:gdelta}
 \WOH=
 \bigcap_{m\geq2}\ \bigcap_{n\geq1}\
 \bigcap_{\mathbf v\in(V\setminus\{0\})^n}
 S_{m,n,\mathbf v}.
\end{equation}
Consequently, $\WOH$ is a $G_\delta$ set in each coding space.
\end{theorem}

\begin{proof}
For a tuple whose entries survive the quotient, Lemma \ref{lem:trace}
runs through exactly all possible restrictions of $f\in B_{X_\mu^*}$ to
the corresponding normalized directions.  Condition
$\mu\in O_{m,n,\mathbf v}$ therefore says exactly that the finite-mesh
criterion \eqref{eq:criterion} holds for that rational-direction tuple.
Tuples with a zero quotient direction impose no condition, as recorded by
$D_{\mathbf v}^{c}$.  The normalized surviving rational directions are
dense in $S_{X_\mu}$.  Proposition \ref{prop:criterion} now proves
\eqref{eq:gdelta}.  All three index sets on the right are countable, and
each $S_{m,n,\mathbf v}$ is $G_\delta$, so their intersection is
$G_\delta$.
\end{proof}

\begin{theorem}[exact complexity]\label{thm:complete}
The class $\WOH$ is $G_\delta$-complete in $\Pinf$ and in $\Bcal$.
\end{theorem}

\begin{proof}
The source paper proves that $\WOH$ is not $F_\sigma$ in either coding space
\cite[Proposition 3.3]{source}.  Theorem \ref{thm:main} gives the
$G_\delta$ upper bound.  The standard Hurewicz completeness theorem for
non-$F_\sigma$ $G_\delta$ subsets of Polish spaces
\cite[Theorem 22.11]{Kechris} therefore makes $\WOH$
$G_\delta$-complete in each space.
\end{proof}

\section{Verification and limitations}

The proof has no computational dependency.  Its quantitative interface is
Lemma \ref{lem:mesh}; all constants have slack by factors of at least two.
The descriptive-set-theoretic interface is finite-dimensional: the compact
cube $[-1,1]^n$ replaces the nonclosed global dual-ball graph in $\Pinf$.
The Hahn--Banach inequalities \eqref{eq:HB} verify that no dual restrictions
are lost or added.

A bounded novelty search on 22 July 2026 covered the run's four cheap
indexes; the exact arXiv id and title; searches for weak octahedrality with
``finite mesh'', ``$G_\delta$'', and ``Borel complexity''; arXiv records and
the current source version; and the cited weak-octahedrality
characterization.  No later paper or preprint stating this exact-complexity
result was found.  Novelty confidence is moderate: the argument is short and
may be independently accessible, but the source explicitly leaves the
complexity open and its stated obstruction is bypassed here.

This packet settles the exact Borel complexity in both coding spaces named by
the source theorem.  It does \emph{not} prove or disprove the single-scale
characterization \eqref{eq:single}.  Human review should focus on the relative
closedness of $\mathfrak C_{\mathbf v}$ when quotient kernels vary, and on the
compact co-projection step defining $O_{m,n,\mathbf v}$.

\begin{thebibliography}{9}
\bibitem{source}
G. L\'opez-P\'erez, E. Mart\'inez Va\~n\'o, and A. Rueda Zoca,
\emph{Computing Borel complexity of some geometrical properties in Banach
spaces}, arXiv:2404.19457 (2024).

\bibitem{HLP}
R. Haller, J. Langemets, and M. P\~oldvere,
\emph{On duality of diameter 2 properties},
J. Convex Anal. \textbf{22} (2015), no.~2, 465--483;
arXiv:1311.2177.

\bibitem{CDDK}
M. C\'uth, M. Dole\v{z}al, M. Doucha, and O. Kurka,
\emph{Polish spaces of Banach spaces: complexity of isometry and isomorphism
classes}, J. Inst. Math. Jussieu \textbf{22} (2023), 1--39.

\bibitem{Kechris}
A. S. Kechris,
\emph{Classical Descriptive Set Theory},
Graduate Texts in Mathematics 156, Springer, 1995.
\end{thebibliography}

\end{document}
